\subsection{Integer-square linearization and central-invertible transport}
\label{sec:complex-low-fourth}
The complex low-product and fourth-moment identities are special cases
of Section~\ref{sec:unified-word-gram}. They retain dual indices even when
the original ring is not self-dual. The following linear consequences
are useful when a complete integral Gram decomposition is too large.
For $P=A_aA_b$, put $q_{ab,v}=\sum_{F(Z)=v}[P:Z]$, let $x_v$ count that
restriction type, and let $B_{ab,v}$ be the coordinatewise restriction
upper bound from the unified theorem. Write $Q_{ab}=[P:P]$.

Introduce $z_{ab,v}=\sum_{F(Z)=v}(N_{A_a,A_b}^Z)^2$ and corresponding
square variables for the known low coefficients. The integer-square inequalities in Section~\ref{sec:unified-word-gram}
apply coefficientwise. Summing them over a restriction type gives
\begin{equation}\label{eq:r4-square-cuts}
 (2k+1)q_{ab,v}-k(k+1)x_v\le z_{ab,v}
          \le B_{ab,v}q_{ab,v}.
\end{equation}
For a low simple, replace $x_v$ by $1$ and $q_{ab,v}$ by its individual
Verlinde multiplicity. Equation~\eqref{eq:r4-fourth-moment} becomes the sum
of these square variables, together with the unit contribution
$\delta_{a,b^*}$. Any finite selection of the lower inequalities is valid;
the implementation supports both all $0\le k<B$ and a stated cutoff.
This gives an exact linear strengthening, with integrality retained in
bounded discovery experiments. Failure of a mixed-integer optimizer to
finish is recorded only as an incomplete search.

\subsubsection{Transport by central invertibles}
Suppose the low frame has exactly one column equal to $e_0$.
A central invertible restricting to the unit must occur in $I(\mathbf1)$,
so the forgetful homomorphism on invertible objects has trivial kernel.
Its image is therefore a subgroup $K$ of the based-ring invertible group,
and each $g\in K$ has exactly one invertible lift.
For any such lift $L$, tensoring by $L$ permutes all central simples and sends
a restriction vector $v$ to the coefficient vector of $gv$. In particular,
the total counts of these restriction types agree, including low objects.

\begin{lemma}[Indicator transport by a central invertible]
\label{lem:central-invertible-indicators}
Let $\mathcal C$ be spherical and let $b\in Z(\mathcal C)$ be
invertible of order $m$, with $\dim b=1$. Write $bX=F(b)\otimes X$.
For every $n\ge1$, $\nu_n(X^*)=\nu_n(X)$; whenever $m\mid n$,
\[
 \nu_n(bX)=\theta_b^n\nu_n(X).
\]
\end{lemma}
\begin{proof}
For the adjoint $I$ of $F$, the projection formula gives
$I(bX)\cong b\otimes I(X)$. For each central simple $Z$, write
$\lambda_Z$ for its monodromy scalar with $b$. The hexagon and
$b^m\cong\one$ imply $\lambda_Z^m=1$, while balancing gives
$\theta_{bZ}=\theta_b\theta_Z\lambda_Z$. Thus
$\theta_{bZ}^n=\theta_b^n\theta_Z^n$ and $\dim(bZ)=\dim Z$.
Apply the spherical center formula \cite[Theorem~4.1]{NS}:
\[
 \nu_n(X)=\frac1{\dim\mathcal C}
 \sum_{Z\in\Irr Z(\mathcal C)}[F(Z):X]\dim Z\,\theta_Z^n.
\]
Relabelling $Z$ by $Z^*$ proves the duality identity, since
$\theta_{Z^*}=\theta_Z$ and $\dim Z^*=\dim Z$; see also
\cite[Remark~4.2]{NS}. No complex conjugation or unitarity is used.
\end{proof}

There is a stronger constraint when $g^2=1$ and $\nu_2(g)=-1$.
The lift has order two and twist $\theta_L=\pm i$. Its monodromy with a
simple $Z$ is a scalar in $\{1,-1\}$: monodromy with invertibles is
multiplicative, and $L^2=\mathbf1$. The balancing identity therefore gives
\[
 \theta_{LZ}=\pm i\theta_Z.
\]
In particular $LZ\not\cong Z$, so this permutation is a fixed-point-free
involution. When the ambient cyclotomic field contains $i$, normalized
relative traces transform by
$(t_j)_j\mapsto((\pm i)^j t_j)_j$.
One can encode the resulting matching with nonnegative integral edge
variables between allowed pairs of restriction/trace types. An edge
joining distinct types contributes one to each marginal; a loop in the
\emph{type} graph contributes two, because it still represents two distinct
simples. This last convention is necessary when relative trace loses the
underlying twist information. These matching equations are necessary only;
for several lifts they do not impose all simultaneous group-action
relations.
